
\appendix

\section{DFT Matrices}\label{app:DFT}

\begin{align*}
	W_2 = e^{-2\pi j / 2} = -1 \\
	A_2 &= 
		\begin{bmatrix}
			1& 1 \\
			1& W_2 \\
		\end{bmatrix}
		&&= 
		\begin{bmatrix}
			1& 1 \\
			1& -1 \\
		\end{bmatrix}	 \\
	W_4 = e^{-2\pi j / 4} = \cos(\pi/2) - j\sin(\pi/2) = -j \\
	A_4 &= 
		\begin{bmatrix}
			1& 1& 1& 1 \\
			1& W_4& W_4^2& W_4^3 \\
			1& W_4^2& W_4^4& W_4^6 \\
			1& W_4^3& W_4^6& W_4^9 \\
		\end{bmatrix}
		&&=
		\begin{bmatrix}
			1& 1& 1& 1 \\
			1& -j& -1& j \\
			1& -1& 1& -1 \\
			1& j& -1& -j \\
		\end{bmatrix} \\
	W_8 = e^{-2\pi j / 8} = \cos(\pi/4) - j\sin(\pi/4) = \frac{1-j}{\sqrt{2}} \\
	A_8 &= 
		\begin{bmatrix}
			1 & 1 & 1 & \dots & 1 \\
			1 & W_8 & W_8^2 & \dots & W_8^7 \\
			1 & W_8^2 & W_8^4 & \dots & W_8^{14} \\
			1 & W_8^3 & W_8^6 & \dots & W_8^{21} \\
			\vdots & \vdots & \vdots & \ddots & \vdots \\
			1 & W_8^7 & W_8^{14} & \vdots & W_8^{49} \\
		\end{bmatrix}
		%TODO wer lustich ist darf das gerne noch vervollständigen
\end{align*}

\[
\begin{aligned}
W_8^0&=1\\
W_8^1&=\frac{1-j}{\sqrt2}\\
W_8^2&=-j\\
W_8^3&=-\frac{1+j}{\sqrt2}\\
W_8^4&=-1\\
W_8^5&=-\frac{1-j}{\sqrt2}\\
W_8^6&=j\\
W_8^7&=\frac{1+j}{\sqrt2}
\end{aligned}
\]
%%%%%%%%%%%%%%%%%%%%%%%%%%%%%%%%%%%%%%%%%%%%%%%%%%%%%%%%%%%%
\section{Common Pole / Zero Forms}

The poles are obtained from the denominator

\[
D(z)=0,
\]

and the zeros from the numerator

\[
N(z)=0.
\]

For second-order polynomials the following standard forms occur frequently.

\begin{center}
\renewcommand{\arraystretch}{1.4}
\begin{tabular}{|p{5.4cm}|p{4.8cm}|p{5.8cm}|}
\hline
\textbf{Polynomial} &
\textbf{Pole / Zero Location} &
\textbf{Comments}
\\ \hline

$1-az^{-1}$

&
$z=a$

&
Real pole/zero.
\\ \hline

$1-2\cos(\omega_0)z^{-1}+z^{-2}$

&
$z=e^{\pm j\omega_0}$

&
Complex-conjugate zeros on the unit circle
($|z|=1$).
\\ \hline

$1-2r\cos(\omega_0)z^{-1}+r^2z^{-2}$

&
$z=re^{\pm j\omega_0}$

&
Complex-conjugate poles/zeros with radius $r$.
\\ \hline

$1-2b\cos(\omega_0)z^{-1}
+(2b-1)z^{-2}$

&
Solve using the quadratic formula

&
Appears in notch and peak filters.
\\ \hline

\end{tabular}
\end{center}

%%%%%%%%%%%%%%%%%%%%%%%%%%%%%%%%%%%%%%%%%%%%%%%%%%%%%%%%%%%%

\paragraph{Useful Identities}

Euler's identity

\[
e^{j\omega}
=
\cos(\omega)+j\sin(\omega)
\]

implies

\[
e^{-j\omega}
=
\cos(\omega)-j\sin(\omega).
\]

Therefore

\[
\boxed{
(z-e^{j\omega_0})
(z-e^{-j\omega_0})
=
z^2-2\cos(\omega_0)z+1
}
\]

or equivalently

\[
\boxed{
1-2\cos(\omega_0)z^{-1}+z^{-2}
=
(1-e^{j\omega_0}z^{-1})
(1-e^{-j\omega_0}z^{-1})
}
\]

Similarly,

\[
\boxed{
1-2r\cos(\omega_0)z^{-1}+r^2z^{-2}
=
(1-re^{j\omega_0}z^{-1})
(1-re^{-j\omega_0}z^{-1})
}
\]

Hence

\[
|z|=
\begin{cases}
1,&\text{unit circle}\\
r,&\text{radius }r
\end{cases}
\]

and

\[
\angle z=\pm\omega_0.
\]

\paragraph{Recognition Rules}

\[
1-2\cos(\omega_0)z^{-1}+z^{-2}
\]

↓

- radius \(=1\)

- angle \(=\pm\omega_0\)

---

\[
1-2r\cos(\omega_0)z^{-1}+r^2z^{-2}
\]

↓

- radius \(=r\)

- angle \(=\pm\omega_0\)

---

\[
1-az^{-1}
\]

↓

- real pole/zero at \(z=a\)
\section{Butterworth polynomials}\label{app:Butterworth}

\begin{tabular}{|l|l|l|l|}
	\hline
	N & K & $\theta_1,\theta_2,\ldots,\theta_K$ & $D(s)$ \\ \hline
	1 & 0 & & $1+s$ \\ \hline
	2 & 1 & $\frac{3\pi}{4}$ & $(1 + 1.4142s + s^2)$ \\ \hline
	3 & 1 & $\frac{4\pi}{6}$ & $(1+s)(1+s+s^2)$ \\ \hline
	4 & 2 & $\frac{5\pi}{8}\:,\:\frac{7\pi}{8}$ & $(1+0.7654s+s^2)(1+1.8478s+s^2)$ \\ \hline
	5 & 2 & $\frac{6\pi}{10}\:,\:\frac{8\pi}{10}$ & $(1+s)(1+0.6180s+s^2)(1+1.6180s+s^2)$ \\ \hline
	6 & 3 & $\frac{7\pi}{12}\:,\:\frac{9\pi}{12}\:,\:\frac{11\pi}{12}$ & $(1+0.5176s+s^2)(1+1.4142s+s^2)(1+1.9319s+s^2)$ \\ \hline
	7 & 3 & $\frac{8\pi}{14}\:,\:\frac{10\pi}{14}\:,\:\frac{12\pi}{14}$ & $(1+s)(1+0.4450s+s^2)(1+1.2470s+s^2)(1+1.8019s+s^2)$ \\ \hline
\end{tabular}
%%%%%%%%%%%%%%%%%%%%%%%%%%%%%%%%%%%%%%%%%%%%%%%%%%%%%%%%%%%%
\section{Kaiser Window}

\[
\alpha=
\begin{cases}
0,&A<21\\[1mm]
0.5842(A-21)^{0.4}
+0.07886(A-21),
&
21\le A\le50\\[1mm]
0.1102(A-8.7),
&
A>50
\end{cases}
\]

\vspace{2mm}

\[
D=
\begin{cases}
0.9222,&A<21\\
\dfrac{A-7.95}{14.36},&A\ge21
\end{cases}
\]

%%%%%%%%%%%%%%%%%%%%%%%%%%%%%%%%%%%%%%%%%%%%%%%%%%%%%%%%%%%%
\section{Useful DSP Relations}

\[
\omega=\frac{2\pi f}{f_s}
\]

\[
\Delta f=\frac{f_s}{N}
\]

\[
f_s'=Lf_s
\]

\[
f_s'=\frac{f_s}{M}
\]

\[
z=\zeta^L
\]

\section{Useful Transform Pairs}
\begin{tabular}{|l|l|}
\hline
Signal & Transform\\
\hline
$\delta(n)$ & $1$\\
\hline
$\delta(n-k)$ & $z^{-k}$\\
\hline
$u(n)$ &
$\dfrac1{1-z^{-1}}$\\
\hline
$a^nu(n)$ &
$\dfrac1{1-az^{-1}}$\\
\hline
$na^nu(n)$ &
$\dfrac{az^{-1}}
{(1-az^{-1})^2}$\\
\hline
\end{tabular}


\section{Idiotenseite}\label{app:Idiotenseite}
\addtocontents{toc}{\protect\setcounter{tocdepth}{-1}}
\input{idiotenseite/trigo/subsections/Winkelargumente}
\input{idiotenseite/trigo/subsections/Quadrantenbeziehungen}
\input{idiotenseite/trigo/subsections/Periodizitaet}
\begin{multicols}{2}
\input{idiotenseite/trigo/subsections/Additionstheoreme}
\input{idiotenseite/trigo/subsections/DoppelHalbwinkel}
\input{idiotenseite/trigo/subsections/Produkte}
\input{idiotenseite/trigo/subsections/SummeDifferenzen}
\subsection{Euler-Formeln}

\paragraph{Euler Identity}

\[
\boxed{
e^{j\varphi}
=
\cos\varphi
+
j\sin\varphi
}
\]

\[
e^{-j\varphi}
=
\cos\varphi
-
j\sin\varphi
\]

\vspace{1mm}

\paragraph{Trigonometrische Darstellung}

\[
\sin\varphi
=
\frac{e^{j\varphi}-e^{-j\varphi}}{2j},
\qquad
\cos\varphi
=
\frac{e^{j\varphi}+e^{-j\varphi}}{2}
\]

\vspace{1mm}

\paragraph{Exponential Rules}

\[
e^{x+jy}
=
e^x\,e^{jy}
=
e^x(\cos y+j\sin y)
\]

\[
e^{j\pi}
=
e^{-j\pi}
=
-1,
\qquad
e^{j\omega}e^{-j\omega}=1
\]

\vspace{1mm}

\paragraph{Complex Numbers}

\[
(a+jb)^*=a-jb
\]

\[
|a+jb|
=
\sqrt{a^2+b^2}
\]

\[
a+jb
=
|z|e^{j\angle z}
\]

\[
\angle(a+jb)
=
\operatorname{atan2}(b,a)
\]

\[
|e^{j\omega}|=1
\]


\end{multicols}
\subsection{ Komplexe Zahlen}

\[
z=a+jb
\]

\[
z^*=a-jb
\]

\[
|z|
=
\sqrt{a^2+b^2}
\]

\[
\angle z
=
\operatorname{atan2}(b,a)
\]

\[
z
=
|z|e^{j\angle z}
\]
\subsection{DFT-Hilfen}

\[
W_N
=
e^{-j\frac{2\pi}{N}}
\]

\[
W_N^N=1
\]

\[
W_N^{k+N}=W_N^k
\]

\[
\sum_{k=0}^{N-1}W_N^{mk}
=
\begin{cases}
N,&m=0\\
0,&m\neq0
\end{cases}
\]
\input{idiotenseite/diverses/subsections/Reihenentwicklung}

